% \subsection{Base Framework} \label{sec:base_framework}

% % A minimal \emph{base framework} built on an \textbf{CPI tree} (rooted clique tree) and a standard peeling routine to compute $(r,s)$-coreness.

% In this section, we present a \emph{base framework} that combines a \textbf{rooted clique tree (CPI)} with a standard peeling routine to compute $(r,s)$-coreness. It avoids explicit $s$-clique enumeration by performing combinatorial counting on CPI paths. The framework applies to all $r\ge 1$; the only variation with $r$ is the deletion primitive—vertex when $r{=}1$, edge when $r{=}2$, and an $r$-clique when $r{\ge}3$.

% % \begin{algorithm}[t]
\small
% \caption{\textsc{$(r,s)$-CND}}
% \label{alg:framework}
% \KwIn{Clique size $s$}
% \KwOut{Core value $\kappa_s(v)$ for every vertex $v$}
% \SetKwFunction{RemoveNode}{RemoveNode}
% \SetKwFunction{CountPreEdges}{$s$-CliqueCount}
% \SetKwProg{Fn}{Function}{:}{}

% $\tree$ $\leftarrow$ build CPI of $G$, pruning paths with $s$\;

% $support \gets$ \textsc{CountPreEdges}$(\mathcal{B}, r, s)$\;

% build min heap $H_{\min}$ using $support$\;

% $c_{min} \gets 0$\;

% \While{$H_{\min}$ is not empty}{
%     $R_{min} \gets \text{ExtractMin}(H_{\min})$\;
%     $c_{min} \gets \max(c_{min}, support(R_{min}))$\;
%     $\kappa_s(R_{min}) \gets c_{min}$\;
%     \ForEach{leaf $\TPath$ intersecting any $rc \in R_{min}$}{
%         $B \gets {\,rc \in R_{min} \mid rc \in \TPath\,}$\;
%         \text{Remove all $r$-cliques in $B$ from $\TPath$}\;
%     }

%     \ForEach{leaf $\TPath$ containing any $rc \in R_{min}$}{
%         Update $support(rc)$ and $H_{\min}$ for all $r$-clique $rc \in \TPath$
%     }

% }

% \Return{$\kappa_s(r)$ for all $r$-cliques $r$ in $\tree$}\;

% \end{algorithm}

\begin{algorithm}[t]
\small
\caption{\textsc{$(r,s)$-\cbnd}}
\label{alg:framework}
\KwIn{Graph $G=(V,E)$, clique sizes $r<s$}
\KwOut{Core value $\kappa_s(R)$ for every $r$-clique $R$}
\SetKwFunction{RemoveNode}{RemoveNode}
\SetKwFunction{CountPreEdges}{$s$-CliqueCount}
\SetKwProg{Fn}{Function}{:}{}

$\tree \leftarrow \textsc{BuildIndex}(G)$\;

\tcp{initialize $s$-support for all $r$-cliques}

$support \gets$ \textsc{$s$-CliqueCount}$(\tree, r, s)$ 

build a min-heap $H_{\min}$ using $support$\;

$c_{min} \gets 0$\;

\While{$H_{\min}$ is not empty}{
    $R_{min} \gets \text{ExtractMin}(H_{\min})$\;
    $c_{min} \gets \max(c_{min}, support(R_{min}))$\;
    $\kappa_s(R_{min}) \gets c_{min}$\;
    
    % remove $R_{min}$ from index
    % Remove all $r$-cliques in $R_{min}$ from $\tree$\;
    \textsc{UpdateIndex}($R_{min}$, $\tree$) \tcp*{remove $R_{min}$ from $\tree$}

    \tcp{update support of affected $r$-cliques}

    \textsc{UpdateSupport}($R_{min}$, $\tree$, $support$, $H_{\min}$)

}

\Return{$\kappa_s(R)$ for all $r$-cliques $R$ in $\tree$}\;

\end{algorithm}


% \paragraph{Framework description.}
% The framework first builds an CPI over $G$ using a degeneracy ordering and pivoted expansion. It then computes the initial counts $\delta(R)$ for every $r$-clique $R$ via Algorithm~\ref{alg:local_s_clique_count}, which evaluates the contribution of each clique path using Lemma~\ref{lem:clique_count} without enumerating $s$-cliques. Next, we initialize a min-heap keyed by $\delta(\cdot)$ and execute a standard peeling loop: repeatedly extract the current minimum $c$, assign $\kappa_s(R){=}c$ to all $r$-cliques with $\delta(R){=}c$, remove these $r$-cliques from the CPI, and decrement affected neighbors using path-level batch updates. The loop terminates when no active $r$-clique remains, and the algorithm outputs $\kappa_s(\cdot)$ for all $r$-cliques.

% % 从

% \paragraph{Complexity analysis.}
% The time complexity of the framework is $O(|E| + |K_r| \log |K_r| + T_{count} + T_{update})$, where $|E|$ is the number of edges in the graph, $|K_r|$ is the number of $r$-cliques, $T_{count}$ is the time taken to count the number of $s$-cliques containing each $r$-clique, and $T_{update}$ is the time taken to update the counts during the peeling process. The space complexity is $O(|E| + |K_r|)$, which includes the space for storing the graph and the CPI.

% % 该算法框架及其简单, 但是在第11行中, 需要从tree中删除$r$-clique, 这一步会很有挑战

% \paragraph{Challenge (line 11).}
% The framework is simple; however, the operation at line~11—removing $r$-cliques from the tree—is non-trivial, mainly because we must maintain the path semantics Lemma~\ref{lem:clique_count} after deletions.

% % 我们在后续的三个小节中, 分别介绍$r=1$, $r=2$和$r\ge3$的情况, 以及如何高效地实现第11行的删除操作. 这三种情况分别对应于顶点删除, 边删除和$r$-clique删除.
% To address this challenge, we proceed in a step-by-step manner in the next three subsections, starting with vertex removal (when $r{=}1$), followed by edge removal (when $r{=}2$), and finally $r$-clique removal (when $r{\ge}3$).

% % In each case, we employ the same method, differing only in the object being removed (vertex when $r{=}1$, edge when $r{=}2$, and an $r$-clique when $r{\ge}3$), we unfold the same method across three settings while keeping one invariant, the same data structures, and the same peeling schedule.
